\documentclass[conference]{IEEEtran}
\usepackage{amsmath, amssymb}
\usepackage{booktabs}
\usepackage{hyperref}
\usepackage{xcolor}

\title{HW1: Introduction to Reinforcement Learning}
\author{Taha Majlesi}

\begin{document}
\maketitle

\begin{abstract}
This report summarizes the HW1 exercises on value-based and policy-based reinforcement learning, including MountainCar control, environment visualization, and a custom Pygame task. We document the experimental configuration, reproducibility practices (seed setting, device logging, and run metadata), and the observed behavior of PPO baselines.
\end{abstract}

\section{Introduction}
The assignment explores foundational RL workflows: environment inspection with Gymnasium, visualization, baseline training with PPO from Stable-Baselines3, and custom environment design. The goal is to align the notebook implementation with the course description while maintaining reproducibility and clean modular structure.

\section{Methods}
\subsection{Environment}
MountainCar-v0 is used as the primary benchmark with a video recorder for episodic rollouts. Optional headless rendering relies on \texttt{xvfb} for Linux; macOS/Windows users skip the system packages.

\subsection{Algorithm}
We employ PPO with Stable-Baselines3. The base configuration mirrors the notebook constants:
\begin{table}[h]
    \centering
    \caption{PPO Hyperparameters (BASE\_CONFIG)}
    \begin{tabular}{ll}
        \toprule
        Parameter & Value \\
        \midrule
        Total timesteps & 20{,}000 \\
        Learning rate & $3\times10^{-4}$ \\
        Discount ($\gamma$) & 0.99 \\
        Rollout steps & 2048 \\
        Batch size & 64 \\
        Video interval & 250 episodes \\
        Seed & 42 \\
        \bottomrule
    \end{tabular}
\end{table}

\section{Reproducibility}
The first notebook cell centralizes imports and seeds \texttt{random}, NumPy, and Torch (including CUDA branches with deterministic flags). Run metadata is persisted to \texttt{run\_info.json} and records commit hash, hardware, hyperparameters, and device availability. A requirements file pins the versions, and setup instructions recommend a virtual environment.

\section{Experiments}
Training PPO for 20k timesteps produces a competent policy on MountainCar. The notebook retains visualization helpers (matplotlib inline rendering and optional video recording) to qualitatively assess progress. For the bonus Pygame grid-world, PPO/DQN can be swapped with identical configuration hooks.

\section{Conclusion}
The HW1 notebook now satisfies the stated structure: setup/configuration precedes environment exploration, followed by model training and qualitative evaluation. Logging and seeding ensure repeatable results, while the LaTeX report documents the configuration expected during grading.

\end{document}












